\documentclass[12pt]{article}
\usepackage[utf8]{inputenc}
\usepackage[T1]{fontenc}
\usepackage{hyperref}
\title{Tunica Language Resources for K-12 Curriculum Development}
\author{Compiled by Claude Sonnet 5 (claude-sonnet-5) for Andr\'as Kornai}
\date{}
\begin{document}
\maketitle

\section{Introduction}
Tunica (ISO 639-3 tun) is a language isolate of Louisiana, historically spoken by the Tunica-Biloxi Tribe. The last first-language speaker died in the twentieth century. Revitalization has been under way since 2010 and is run by the tribe with Tulane University. One source reports that Ethnologue changed the language's status from ``dormant'' to ``reawakened''. This report lists only what the search found.

\section{Primary Educational Resources}
\begin{itemize}
\item The Tunica-Biloxi Tribe's Language \& Culture Revitalization Program runs four 8-week Tunica class series a year, in fall and spring. It also develops teaching materials and operates the CERC library. \url{https://tunicabiloxi.org/tribal-info/departments/language-culture-revitalization-program}
\item The program has run a summer language camp since 2011. Children there learned greetings, numbers, colors and basic verbs. See the grant news item below.
\item Announcement that the tribe hired a linguist to develop resources and lesson plans. \url{https://www.tunicabiloxi.org/the-tunica-biloxi-tribe-of-louisiana-welcomes-new-linguist-to-language-and-culture-revitalization-program-staff/}
\end{itemize}

\section{Linguistic and Reference Materials}
\begin{itemize}
\item Tunica Language Project, Tulane University: about the project. \url{https://tunica.wp.tulane.edu/about/about-the-tunica-language-project/}
\item Tunica Language Project, Tulane (further pages). \url{https://tunica.wp.tulane.edu/?p=37} and \url{https://tunica.wp.tulane.edu/?p=53}
\item EBSCO Research Starters, ``Tunica language''. \url{https://www.ebsco.com/research-starters/language-and-linguistics/tunica-language}
\end{itemize}
The search did not surface Mary Haas's Tunica texts or grammar. Do not infer that they are online.

\section{Religious and Community Resources}
\begin{itemize}
\item Tunica-Biloxi Tribe of Louisiana, official website. \url{https://www.tunicabiloxi.org/?p=779}
\item Tribe news: grant for language preservation. It describes a \$748,200 Administration for Native Americans grant that funds five adult apprentices, who are expected to start language nests afterward. \url{https://www.tunicabiloxi.org/tunica-biloxi-tribe-receives-grant-for-language-preservation/}
\item Other tribal pages. \url{https://tunicabiloxi.org/?p=242} and \url{https://www.tunicabiloxi.org/?p=797}
\end{itemize}

\section{Digital and Audio Resources}
No audio, video or app resources were identified in this search. The Tulane project site may host materials; its contents were not examined.

\section{Curriculum Development Context}
The tribe already runs a language department with classes, camps and funded apprentices, so any school curriculum should be built with the Language \& Culture Revitalization Program as lead partner. The sources found date mostly from 2018 or earlier. Current staffing, grant status and class schedules need confirming directly with the tribe.

\end{document}